\documentclass[10pt]{article}

\usepackage[utf8]{inputenc}

\usepackage[T1]{fontenc}

\usepackage{amsmath}

\usepackage{amsfonts}

\usepackage{amssymb}

\usepackage[version=4]{mhchem}

\usepackage{stmaryrd}

\usepackage{graphicx}

\usepackage[export]{adjustbox}

\graphicspath{ {./images/} }

\usepackage{bbold}



\begin{document}

БETE-СОЛПИТЕРА УРАВНЕНИЕ - релятивистское соотношение для двухчастичной Грина функции $D\left(x_{1}, x_{2} ; x_{1}^{\prime}, x_{2}^{\prime}\right)$ системы двух частиц (или полей):



\[

\begin{gathered}

D\left(x_{1}, x_{2} ; x_{1}^{\prime}, x_{2}^{\prime}\right)=\bar{D}\left(x_{1}, x_{2} ; x_{1}^{\prime}, x_{2}^{\prime}\right)+ \\

+\int K\left(x_{1}, x_{2} ; x_{3}, x_{4}\right) D\left(x_{3}, x_{4} ; x_{1}^{\prime}, x_{2}^{\prime}\right) d^{4} x_{3} d^{4} x_{4}

\end{gathered}

\]



$\left(x_{1}, x_{2}, x_{1}^{\prime}, x_{2}^{\prime}\right.$ - начальные и конечные 4-мерные координаты частиц). Сформулировано Х. Бете (Н. Веthe) и Э. Солпитером (E. Salpeter) в 1951 для описания связанных состояний системы частиц 1 и 2, к-рым отвечают полюсы функции $D$ [в этом случае в уравнении $\left(^{*}\right)$ отсутствует неоднородный член $\vec{D}$, не содержащий этих полюсов], и опирается на инвариантную теорию возмущений в форме Фейнмана диаграмм. Б.-С.у. связывает полную функцию Грина двух частиц $D\left(x_{1}, x_{2} ; x_{1}^{\prime}, x_{2}^{\prime}\right)$, понимаемую как сумма всех диаграмм Фейнмана (рис. 1, левая часть), с определенной, топологически выделенной частью этой суммы $\bar{D}\left(\mathrm{x}_{1}\right.$, $\left.x_{2} ; x_{1}^{\prime}, x_{2}^{\prime}\right)$ (рис. 1 , первое слагаемое правой части), представляющей собой сумму всех двухчастично неприводимых диаграмм в $\mathrm{t}$-канале [в к-ром кинематич. переменная $t=\left(p_{1}+p_{2}\right)^{2}$, где $p_{1}, p_{2}-4$-импульсы частиц 1,2$]$, то есть таких диаграмм, к-рые нельзя разбить на две связные части, содержашие точки $x_{1}, x_{2}$ и $x_{1}^{\prime}, x_{2}^{\prime}$, разорвав только две линии, идущие в направлении $t$-канала.



Ядро $K$ Б.-С. у. явным образом выражается через сумму двухчастично неприводимых диаграмм $\bar{D}$ и функции Грина свободных частиц [второе слагаемое в правой части рис. 1



\begin{center}

\includegraphics[max width=\textwidth]{2023_07_08_5f40f101bfc11d6c5122g-1(1)}

\end{center}



Рис. 1.



отвечает интегральному члену в уравнении $\left.\left({ }^{*}\right)\right]$. Оно строится на основе лагранжиана взаимодействия частиц с полем, но само поле в уравнение $\left({ }^{*}\right)$ явно не входит. Поскольку, кроме теории возмущений, других конструктивных методов вычисления ядра (так же, как и неоднородного члена) точного Б.-С. у. не существует, его следует рассматривать только как удобное соотношение, позволяющее неявным образом выразить всю сумму диаграмм Фейнмана через их двухчастично неприводимую часть.



Часто под Б.-С.у. понимают приближенное уравнение (так наз. лестничное прибли жен ие), к-рое получается, если ограничиться в сумме двухчастично неприводимых диаграмм низшим порядком теории возмущений, то



\includegraphics[max width=\textwidth, center]{2023_07_08_5f40f101bfc11d6c5122g-1}

есть однократным обменом квантом поля между двумя взаимодействующими частицами (рис. 2). В этом приближении Б.-С. у. обычно используется для релятивистского описания связанных состояний системы двух слабо взаимодействующих частиц. В нерелятивистском пределе Б.-С.. перестает зависеть от двух различных времен и переходит в уравнение Шредингера с соответствующим потенциалом.



Уравнение $\left({ }^{*}\right)$ можно понимать и как уравнение непосредственно для амплитуды рассеяния двух частиц. В этом случае $D$ следует считать амплитудой, а $\bar{D}-$ ее неприводимой частью. Упомянутое выше соотношение между функциями $\bar{D}$ и $K$ сохраняется. Учитывая перекрестную симметрию



\section{2 БЕТЕ-СОЛПИТЕРА}

амплитуды рассеяния, связывающую различные каналы реакции (этому свойству удовлетворяют диаграммы Фейнмана во всех порядках теории возмущений), можно использовать Б.-С. у. для описания взаимодействия часдиц в $s$-канале [в к-ром кинематич. переменная $s=\left(p_{1}-p_{1}^{\prime}\right)^{2}$; см. рис. 1], то есть рассматривать с его помощью столкновение частицы 1 с античастицей $\tilde{l}^{\prime}$ с превращением их в частицу 2 и античастицу $\tilde{2}^{\prime}$. Такая трактовка Б.-С.у. положена в основу мультипериферич. модели процессов множественного рождения частиц при высоких энергиях. В этом случае из Б.-С. у. удается получить аналогичное, но более простое уравнение для мнимой части амплитуды рассеяния частиц $l$ и $\widetilde{l}^{\prime}$ в $s$-канале, к-рая посредством оптической теоремы связана с полным сечением упомянутых процессов.



Jum.: [1] S a lpeter E. E., B e the H. A., A relativistic equation for bound-state problems, «Phys. Rev.», 1951, v. 84, p. 1232; [2] Ш ве бе р С., Введение в релятивистскую квантовую теорию поля, пер. с англ., М., 1963, гл. 17, §6; [3] Дунаевский А. М., Ройзе н И. И., О виковском повороте в рамках уравнения Бете - Солпитера, «Ядерная физика», 1971, т. 14, с. 855; [4] Л и фши ц Е. М., Питаевский Л.П., Релятивистская квантовая теория, ч. 2, M., 1971. И. И. Ройзен.



БИАНКИ ТОЖДЕСТВО для тензора кривизны $R_{j, k l}^{i}$ симметричной связности, согласованной с метрикой,- тождество



\[

\nabla_{m} R_{i, k l}^{n}+\nabla_{l} R_{i, m k}^{n}+\nabla_{k} R_{i, l m}^{n}=0 .

\]



(см. [1], [2]). Аналогичное равенство имеет место для связностей в расслоениях, а именно:



(cM. [3]).



\[

D_{[\gamma} R_{\alpha \beta]}=0

\]



См. также Кривизны форма.



Jum.: [1] B i a n c h i L., Lezioni di geometria differenziale, v. 1, 3 ed., Bologna, 1923; [2] Р аше в с к и й П.К., Риманова геометрия и тензорный анализ, М., 1967; [3] Т роф и м ов В. В., Введение в геометрию многообразий с симметриями, М., $1989 . \quad$ В. В. Трофимов. БИБЕРБАХА ГИПОТЕЗА - см. ОӘнолистности условия. БИГАРМОНИЧЕСКАЯ ФУНКЦИЯ - ФУНКЦИЯ



\[

u(x)=u\left(x_{1}, \ldots, x_{n}\right)

\]



действительных переменных, определенная в области $D$ евклидова пространства $\mathbb{R}^{n}, n \geqslant 2$, имеющая непрерывные частные производные до 4-го порядка включительно и удовлетворяющая в $D$ уравнению



\[

\Delta^{2} u \equiv \Delta(\Delta u)=0,

\]



где $\Delta$ - оператор Лапласа. Это уравнение называется бигармоническим уравнением. Класс Б. Ф. включает Класс гармонических функций и является подклассом класса полигармонииеских функций. Каждая Б. ф. есть аналитич. функция от координат $x_{i}$.



Наибольшее значение с точки зрения приложений имеют Б. ф. $u\left(x_{1}, x_{2}\right)$ двух переменных. Такие Б. ф. представимы при помощи гармонич. функций $u_{1}, u_{2}$ или $v_{1}, v_{2}$ в виде



\[

u\left(x_{1}, x_{2}\right)=x_{1} u_{1}\left(x_{1}, x_{2}\right)+u_{2}\left(x_{1}, x_{2}\right)

\]



или



\[

u\left(x_{1}, x_{2}\right)=\left(r^{2}-r_{0}^{2}\right) v_{1}\left(x_{1}, x_{2}\right)+v_{2}\left(x_{1}, x_{2}\right),

\]



где $r^{2}=x_{1}^{2}+x_{2}^{2}$, а $r_{0}^{2}-$ постоянная. Основная краевая задача для Б. ф. состоит в следующем: найти Б. ф. в области $D$, непрерывную вместе с производными 1-го порядка в замкнутой области $\bar{D}=D \cup C$ и удовлетворяющую на границе $C$ условиям



\[

\left.u\right|_{C}=f_{1}(s),\left.\quad \frac{\partial u}{\partial n}\right|_{C}=f_{2}(s),

\]



где $\partial u / \partial n-$ производная по нормали к $C$, а $f_{1}(s), f_{2}(s)$ заданные непрерывные функции дуги $s$ на контуре $C$. Указанные выше представления Б. ф. позволяют получить решение задачи $\left(\left(^{*}\right)\right.$ в явном виде в случае круга $D$, исходя из интеграла Пуассона для гармонич. функций.





\end{document}